%---------------------------------------
% Packages
%---------------------------------------
% === Document Setup and Language ===
\usepackage[english]{babel}        % English language
\usepackage[utf8]{inputenc}        % Unicode support
\usepackage[T1]{fontenc}           % Font encoding
\usepackage{subfiles}              % Allow for subfiles which can be compiled standalone
\setcounter{secnumdepth}{2}

% === Math Packages ===
\usepackage{amsfonts, amsmath, amssymb, amsthm}  % AMS math packages
\usepackage{dsfont}                % For \mathds{1}
\usepackage{mathtools}             % Enhanced math tools
\usepackage{accents}               % Mathematical accents

% === Typography and Fonts ===
%\usepackage{mathptmx}              % Times New Roman font
\usepackage{newtxtext,newtxmath}   % Enhanced Times New Roman with math support
\usepackage{csquotes}              % Smart quotes and quotation management
\linespread{1.05}                  % Line spacing

% === Graphics and Diagrams ===
\usepackage{tikz-cd, tikz}         % Diagrams and drawings
\usetikzlibrary{
	positioning, 	% For relative positioning (below=of, right=of)
	calc,        	% For coordinate calculations (finding midpoints)
	fit,			% Needed for the bounding box (fit)
	patterns,
	decorations.markings,
	decorations.pathmorphing,	% Needed for squiggly arrows
	backgrounds
}
\usepackage{graphicx}              % Image handling
\usepackage{booktabs}              % \toprule, \midrule, \bottomrule for tables

% === Page Layout ===
\usepackage{geometry}
\geometry{left=1.3in,right=1.2in}
\setlength{\parskip}{\smallskipamount}
\setlength{\parindent}{0pt}

% === Lists and Enumeration ===
\usepackage{enumerate}             % Enhanced enumeration
\renewcommand{\labelenumi}{\normalfont{(\arabic{enumi})}} % Default to [(1)]
\usepackage{etoolbox}              % Programming tools for LaTeX

% === Hyperlinks and References ===
\usepackage[hyphens]{url}
\urlstyle{same}
\usepackage{bookmark}
\usepackage{hyperref}              % Must be loaded after most other packages
\hypersetup{
	colorlinks=true,
	linkcolor=black,
	filecolor=black,
	urlcolor=black,
	citecolor=black,
	bookmarksdepth=subsubsection
}

% === Bibliography ===
\usepackage[
backend=biber,
natbib,
style=alphabetic,
maxcitenames=4,
maxalphanames=4,
maxbibnames=99
]{biblatex}
\renewbibmacro{in:}{}              % Suppresses "In: " before journal names
\addbibresource{C:/Users/LocalAdmin/OneDrive/Documenten OneDrive/A - Regensburg/Research and notes/Main_Bibliography.bib}
%\addbibresource{Bibliography.bib}

% === Optional Settings (uncomment if needed) ===
% \setcounter{secnumdepth}{2}
% \setcounter{tocdepth}{2}

%---------------------------------------
% Gestalten
%---------------------------------------
\DeclareMathOperator{\fpqc}{fpqc}

%---------------------------------------
% Algebraic geometry / motivic cohomology
%---------------------------------------
\DeclareMathOperator{\KGL}{KGL}
\DeclareMathOperator{\Nis}{Nis}
\DeclareMathOperator{\Et}{Et}
\DeclareMathOperator{\cdh}{cdh}
\DeclareMathOperator{\vdim}{vdim}
\DeclareMathOperator{\Krdim}{Kr.dim}
\DeclareMathOperator{\qcqs}{qcqs}
\DeclareMathOperator{\Sm}{Sm}
\DeclareMathOperator{\SSeq}{SSeq}
\DeclareMathOperator{\Sym}{Sym}
\DeclareMathOperator{\KH}{KH}

%---------------------------------------
% Topos theory
%---------------------------------------
\newcommand{\ssquare}{\,\square\,}
\DeclareMathOperator{\Aff}{Aff}
\DeclareMathOperator{\Mono}{Mono}
\DeclareMathOperator{\Gerb}{Gerb}
\DeclareMathOperator{\hyp}{hyp}
\DeclareMathOperator{\mono}{mono}
\DeclareMathOperator{\cons}{cons}
\DeclareMathOperator{\wcons}{wcons}
\DeclareMathOperator{\quottext}{quot}
\DeclareMathOperator{\epi}{epi}
\DeclareMathOperator{\MonoCong}{MonoCong}
\DeclareMathOperator{\EpiCong}{EpiCong}
\DeclareMathOperator{\Cong}{Cong}
\DeclareMathOperator{\HypCong}{HypCong}
\DeclareMathOperator{\Mdl}{Mdl}
\DeclareMathOperator{\Acyc}{Acyc}
\DeclareMathOperator{\Sat}{Sat}
\DeclareMathOperator{\SSat}{SSat}
\DeclareMathOperator{\Conn}{Conn}
\DeclareMathOperator{\EffEpi}{EffEpi}
\DeclareMathOperator{\Epi}{Epi}
\DeclareMathOperator{\Cotr}{Cotr}
\DeclareMathOperator{\All}{All}
\DeclareMathOperator{\Pos}{Pos}
\DeclareMathOperator{\post}{post}
\DeclareMathOperator{\El}{El}
\DeclareMathOperator{\cocompl}{cocompl}
\newcommand{\Topos}{\catname{Topos}}
\newcommand{\Logos}{\catname{Logos}}
\newcommand{\bbCat}{\bbC \mathrm{at}}
\newcommand{\bbLog}{\bbL \mathrm{ogos}}
\newcommand{\bbTop}{\bbT \mathrm{opos}}
%\renewcommand{\top}{\mathrm{top}}
\DeclareMathOperator{\Geom}{Geom}
\DeclareMathOperator{\Pt}{Pt}
\DeclareMathOperator{\im}{im}				% Image
\DeclareMathOperator{\coim}{coim}			% Coimage
\DeclareMathOperator{\Act}{Act}
\DeclareMathOperator{\Bun}{Bun}
\DeclareMathOperator{\Site}{Site}
\DeclareMathOperator{\cont}{cont}
\DeclareMathOperator{\cocont}{cocont}
\DeclareMathOperator{\wccolim}{w.c.colim}

%---------------------------------------
% Miscellaneous
%---------------------------------------
\DeclareMathOperator{\Atl}{Atl}
\DeclareMathOperator{\DiffStk}{DiffStk}
\DeclareMathOperator{\SepStk}{SepStk}
\renewcommand{\H}{\mathrm{H}}
\newcommand{\SH}{\mathrm{SH}}
\DeclareMathOperator{\htp}{htp}
\DeclareMathOperator{\rep}{rep}
\DeclareMathOperator{\Loc}{Loc}
\DeclareMathOperator{\modtext}{mod}
\newcommand{\absmod}[1]{\abs{#1}_{\modtext}}
\DeclareMathOperator{\sep}{sep}
\DeclareMathOperator{\PB}{PB}
\DeclareMathOperator{\PF}{PF}
\DeclareMathOperator{\NerveDiff}{N^{\mathrm{Diff}}}
\newcommand{\Deltaalg}[1]{\Delta^{#1}_{\mathrm{alg}}}
%\newcommand{\Act}[1]{#1 \textup{Act}}
\DeclareMathOperator{\shear}{shear}
\DeclareMathOperator{\Sph}{Sph}
\DeclareMathOperator{\Sub}{Sub}
\DeclareMathOperator{\Open}{Open}
\DeclareMathOperator{\open}{open}
\DeclareMathOperator{\Corr}{Corr}
\DeclareMathOperator{\bbMap}{\mathbb{M} ap}
\DeclareMathOperator{\bbHom}{\mathbb{H} om}
\DeclareMathOperator{\gen}{gen}
\DeclareMathOperator{\Bgl}{B_{gl}}
\newcommand{\PrnBdl}[2]{\textup{PrnBdl}_{#1}(#2)}
\newcommand{\catnameb}[1]{{\normalfont\textbf{#1}}}
\DeclareMathOperator{\bLieGrpd}{\catnameb{LieGrpd}}
\newcommand{\cFun}{\Ff \textup{un}}
\DeclareMathOperator{\Tr}{Tr}
\DeclareMathOperator{\tr}{tr}
\DeclareMathOperator{\Shp}{Shp}
\DeclareMathOperator{\THH}{THH}

%---------------------------------------
% Higher Algebra
%---------------------------------------

\newcommand{\Assoc}{\mathcal{A}\mathit{ssoc}}
\newcommand{\Comm}{\mathcal{C}\!\,\mathit{omm}}
\newcommand{\Lie}{\mathcal{L}\mathit{ie}}
\newcommand{\Triv}{\mathcal{T}\!\!\mathit{riv}}
\newcommand{\oTens}{\mathcal{T}\!\!\mathit{ens}}
\newcommand{\oLMod}{\mathcal{L}\!\mathcal{M}\mathit{od}}
\newcommand{\oRMod}{\mathcal{R}\!\mathcal{M}\mathit{od}}
\newcommand{\oBMod}{\mathcal{B}\!\mathcal{M}\mathit{od}}
\newcommand{\oMod}{\mathcal{M}\mathit{od}}
\newcommand{\oAlg}{\mathcal{A}\mathit{lg}}
\newcommand{\oEnd}{\mathcal{E}\mathit{nd}}
\newcommand{\oCMon}{\mathcal{C}\!\mathcal{M}\mathit{on}}
\newcommand{\oCGrp}{\mathcal{C}\mathcal{G}\mathit{rp}}
\newcommand{\oSp}{\mathcal{S}\mathit{p}}
\newcommand{\Mult}{\mathcal{M}}
\DeclareMathOperator{\oCat}{\mathcal{C}\!\mathit{at}}
\DeclareMathOperator{\oFun}{\mathcal{F}\!\mathit{un}}
%\newcommand{\Mult}{\mathcal{M}\mathit{ult}}
\DeclareMathOperator{\Env}{Env}
\DeclareMathOperator{\Op}{Op}
\DeclareMathOperator{\AdTrip}{AdTrip}
\DeclareMathOperator{\rev}{rev}
\DeclareMathOperator{\tiso}{iso}
\DeclareMathOperator{\oTw}{\overline{\Tw}}
\DeclareMathOperator{\lax}{lax}
\DeclareMathOperator{\NTw}{NTw}
\newcommand{\OpCocart}{\Oo^{\amalg}}
\newcommand{\OpCart}{\Oo^{\times}}
\DeclareMathOperator{\Day}{Day}
\newcommand{\oDay}{\mathcal{D}\mathit{ay}}
\newcommand{\dlax}{\textup{-lax}}
\newcommand{\textred}{\textup{red}}
\newcommand{\exc}{\textup{exc}}
\DeclareMathOperator{\BiFun}{BiFun}
\DeclareMathOperator{\BV}{BV}
\renewcommand{\Bar}{\mathrm{Bar}}
\newcommand{\dnil}{\textup{-nil}}
\DeclareMathOperator{\Ext}{Ext}
\DeclareMathOperator{\Ex}{Ex}
\DeclareMathOperator{\Tor}{Tor}
\newcommand{\enr}{\textup{enr}}
\newcommand{\heart}{\heartsuit}
\newcommand{\dsm}{\textup{-sm}}
\DeclareMathOperator{\Free}{Free}

%---------------------------------------
% Stable Homotopy Theory
%---------------------------------------
\DeclareMathOperator{\ku}{ku}
\DeclareMathOperator{\KU}{KU}
\DeclareMathOperator{\MU}{MU}
\DeclareMathOperator{\MO}{MO}
\DeclareMathOperator{\MSO}{MSO}
\DeclareMathOperator{\BO}{BO}
\DeclareMathOperator{\bGr}{\mathbf{Gr}}
\DeclareMathOperator{\BUP}{BUP}
\DeclareMathOperator{\bBUP}{\mathbf{BUP}}
\DeclareMathOperator{\Th}{Th}
\DeclareMathOperator{\Pic}{Pic}
\DeclareMathOperator{\pic}{pic}
\DeclareMathOperator{\pre}{pre}
\DeclareMathOperator{\Line}{Line}
%\DeclareMathOperator{\Triv}{Triv}
\DeclareMathOperator{\PT}{PT}
\newcommand{\interiorsymb}[1]{ {\kern0pt#1}^{\mathrm{o}}}
\newcommand{\PSp}{\catname{PSp}}
\DeclareMathOperator{\Emb}{Emb}
\newcommand{\Ntop}{N^{\textup{top}}}
\newcommand{\Homtop}{\Hom^{\textup{top}}}
\DeclareMathOperator{\fw}{fw}
\DeclareMathOperator{\Shift}{Shift}
\DeclareMathOperator{\Der}{Der}

%---------------------------------------
% Category Theory Operators
%---------------------------------------

% Maps and functors
\DeclareMathOperator{\Hom}{Hom}
\DeclareMathOperator{\Aut}{Aut}
\DeclareMathOperator{\End}{End}
\DeclareMathOperator{\iHom}{\underline{\operatorname{Hom}}}			% Internal Hom
\DeclareMathOperator{\ihom}{\underline{\operatorname{hom}}}			% Internal Hom
\DeclareMathOperator{\uHom}{\underline{\operatorname{Hom}}}			% Internal Hom
\DeclareMathOperator{\Map}{Map}
\DeclareMathOperator{\map}{map}
\DeclareMathOperator{\Path}{Path}
\DeclareMathOperator{\Fun}{Fun}
\newcommand{\FunL}{\Fun^{\textup{L}}}
\newcommand{\FunR}{\Fun^{\textup R}}
\newcommand{\ulFun}{\ul{\Fun}}
\DeclareMathOperator{\Nat}{Nat}

% Fibrations and Related Concepts
\DeclareMathOperator{\Lift}{Lift}
\DeclareMathOperator{\lift}{lift}
\DeclareMathOperator{\cocart}{cocart}
\DeclareMathOperator{\CoCart}{CoCart}
\DeclareMathOperator{\Cocart}{Cocart}
\DeclareMathOperator{\isCocart}{\texttt{isCocart}}
\DeclareMathOperator{\Cart}{Cart}
\DeclareMathOperator{\RFib}{RFib}
\DeclareMathOperator{\LFib}{LFib}
\DeclareMathOperator{\St}{St}
\DeclareMathOperator{\Un}{Un}
\DeclareMathOperator{\cart}{cart}
\newcommand{\dcart}{\textup{-cart}}
\DeclareMathOperator{\cc}{cc}
\DeclareMathOperator{\ct}{ct}
\DeclareMathOperator{\Str}{Str}
\DeclareMathOperator{\el}{el}

%---------------------------------------
% Algebraic Geometry
%---------------------------------------
\DeclareMathOperator{\Spec}{Spec}
\DeclareMathOperator{\QCoh}{QCoh}
\DeclareMathOperator{\IndCoh}{IndCoh}
\DeclareMathOperator{\Perf}{Perf}
\DeclareMathOperator{\Proj}{Proj}
\DeclareMathOperator{\Zar}{Zar}
\DeclareMathOperator{\Sch}{Sch}
\newcommand{\ad}{\mathrm{ad}}
\newcommand{\corp}{\mathrm{corp}}
\newcommand{\loc}{\mathrm{loc}}
\newcommand{\mot}{\mathrm{mot}}
\newcommand{\ultra}{\mathrm{ultra}}

%---------------------------------------
% Other math-operators
%---------------------------------------
\DeclareMathOperator{\Ar}{Ar}
\DeclareMathOperator{\BC}{BC}
\DeclareMathOperator{\GL}{GL}
\DeclareMathOperator{\Gr}{Gr}
\DeclareMathOperator{\Ho}{Ho}
\let\Im\relax
\DeclareMathOperator{\Im}{Im}
\DeclareMathOperator{\Ind}{Ind}
\DeclareMathOperator{\Iso}{Iso}
\DeclareMathOperator{\LConst}{LConst}
\DeclareMathOperator{\Nm}{Nm}
\DeclareMathOperator{\Pro}{Pro}
\DeclareMathOperator{\Lan}{Lan}
\DeclareMathOperator{\Ran}{Ran}
\DeclareMathOperator{\Rex}{Rex}
\DeclareMathOperator{\PSh}{PSh}
\DeclareMathOperator{\Shape}{Shape}
\DeclareMathOperator{\Shv}{Shv}
\usepackage{relsize}
\newcommand{\simp}{\mathlarger{\boldsymbol{\Delta}}}
\newcommand{\bOmega}{{\boldsymbol{\Omega}}}
\DeclareMathOperator{\Span}{Span}
\newcommand{\SpanTwo}{{\textup{\textsc{Span}}}_2}
\newcommand{\SpanTwoP}[4]{\SpanTwo(#1,#2)_{#4,#3}}
\DeclareMathOperator{\NSpan}{NSpan}
\DeclareMathOperator{\Tw}{Tw}
\DeclareMathOperator{\LM}{LM}
\let\Re\relax
\DeclareMathOperator{\Re}{Re}
\DeclareMathOperator{\Fact}{Fact}

\DeclareMathOperator{\SD}{SD}				% Suave dual
\DeclareMathOperator{\PD}{PD}				% Prim dual
\DeclareMathOperator{\Suave}{Suave}			% Suave objects
\DeclareMathOperator{\Prim}{Prim}			% Prim objects

\DeclareMathOperator{\shift}{shift}
\DeclareMathOperator{\act}{act}
\DeclareMathOperator{\supp}{supp}
\DeclareMathOperator{\ex}{ex}
\DeclareMathOperator{\lex}{lex}
\DeclareMathOperator{\rex}{rex}
\DeclareMathOperator{\chain}{chain}
\DeclareMathOperator{\sd}{sd}
\DeclareMathOperator{\qc}{qc}
\DeclareMathOperator{\Poly}{Poly}
\DeclareMathOperator{\Mot}{Mot}

%---------------------------------------
% Abbreviations
%---------------------------------------

\newcommand{\can}{\mathrm{can}}             % Canonical
\newcommand{\ho}{\mathrm{ho}}               % Homotopy category
\newcommand{\nondeg}{\mathrm{nondeg}}       % Non-degenerate
\newcommand{\univ}{\mathrm{univ}}           % Universal
\newcommand{\prop}{\textup{pr}}
\newcommand{\inv}{\textup{inv}}
\DeclareMathOperator{\ac}{ac}				% Associated category
\DeclareMathOperator{\pb}{pb}				% Pullback

\DeclareMathOperator{\fgt}{fgt}				% Forgetful functor
\DeclareMathOperator{\id}{id}				% Identity
\DeclareMathOperator{\pr}{pr}				% Projection
\DeclareMathOperator{\incl}{incl}				% Inclusion
\DeclareMathOperator{\const}{const}			% Forgetful
\DeclareMathOperator{\sk}{sk}				% Skeleton
\DeclareMathOperator{\cosk}{cosk}			% Coskeleton
\DeclareMathOperator{\gl}{gl}               % Global
\DeclareMathOperator{\grp}{grp}             % Group completion
\DeclareMathOperator{\qis}{qis}             % Quasi-isomorphism
\DeclareMathOperator{\cyl}{cyl}             % Cylinder
\DeclareMathOperator{\surj}{surj}             % Surjective
\newcommand{\et}{\textup{ét}}				% Étale
\newcommand{\coh}{\mathrm{coh}}				% Coherent
\newcommand{\acc}{\mathrm{acc}}				% Accessible
\renewcommand{\small}{\mathrm{small}}

\DeclareMathOperator{\free}{free}
\DeclareMathOperator{\fin}{fin}
\DeclareMathOperator{\pt}{pt}
\DeclareMathOperator{\aug}{{aug}}
\DeclareMathOperator{\st}{st}

\DeclareMathOperator{\cofib}{cofib}
\DeclareMathOperator{\fib}{fib}
\DeclareMathOperator{\colim}{colim}
\let\lim\relax
\DeclareMathOperator{\lim}{lim}
\DeclareMathOperator{\coeq}{coeq}
\DeclareMathOperator{\coker}{coker}
\newcommand{\ev}{\textup{ev}}
\newcommand{\coev}{\textup{coev}}
\DeclareMathOperator{\res}{res}
\DeclareMathOperator{\ind}{ind}
\DeclareMathOperator{\coind}{coind}

%---------------------------------------
% Categories
%---------------------------------------
\newcommand{\catname}[1]{{\textup{#1}}}
\newcommand{\Set}{\catname{Set}}
\newcommand{\An}{\catname{An}}
\newcommand{\CondAn}{\catname{CondAn}}
\newcommand{\Ab}{\catname{Ab}}
\newcommand{\Alg}{\catname{Alg}}
\newcommand{\CAlg}{\catname{CAlg}}
\newcommand{\Cat}{\catname{Cat}}
\newcommand{\CAT}{\catname{CAT}}
\newcommand{\CRing}{\catname{CRing}}
\newcommand{\qCat}{\catname{qCat}}
\newcommand{\CGrp}{\catname{CGrp}}
\newcommand{\CMon}{\catname{CMon}}
\newcommand{\Fin}{\catname{Fin}}
\newcommand{\FinGrpd}{\catname{FinGrpd}}
\newcommand{\Grp}{\catname{Grp}}
\newcommand{\Grpd}{\catname{Grpd}}
\newcommand{\Mod}{\catname{Mod}}
\newcommand{\coMod}{\catname{coMod}}
\newcommand{\Mon}{\catname{Mon}}
\newcommand{\LMod}{\catname{LMod}}
\newcommand{\RMod}{\catname{RMod}}
\newcommand{\BMod}{\catname{BMod}}
\newcommand{\Rep}{\catname{Rep}}
\renewcommand{\Top}{\catname{Top}}
\newcommand{\TopGrpd}{\catname{TopGrpd}}
\newcommand{\bTopGrpd}{\mathbf{TopGrpd}}
\newcommand{\Spc}{\catname{Spc}}
\newcommand{\Sp}{\catname{Sp}}
\newcommand{\Vect}{\catname{Vect}}
\newcommand{\Orb}{\textup{Orb}}
\newcommand{\Orbit}{\textup{Orbit}}
\newcommand{\Glo}{\textup{Glo}}

\newcommand{\sAn}{\s\An}
\newcommand{\dAn}{\mathrm{d}\An}
\newcommand{\bVect}{\mathbf{V}\catname{ect}}
\newcommand{\LCHaus}{\catname{LCHaus}}
\newcommand{\CHaus}{\catname{CHaus}}
\newcommand{\ProFin}{\catname{ProFin}}
\newcommand{\AnRing}{\catname{AnRing}}
\newcommand{\AnStck}{\catname{AnStck}}
\newcommand{\Ring}{\catname{Ring}}
\newcommand{\twoCAlg}{\catname{2CAlg}}
\newcommand{\Idem}{\catname{Idem}}
\newcommand{\coIdem}{\catname{coIdem}}
\newcommand{\StRing}{\catname{StRing}}
\newcommand{\PreStRing}{\catname{PreStRing}}
\newcommand{\Gest}{\catname{Gest}}
\newcommand{\RingGest}{\catname{RingGest}}
\newcommand{\LieGrpd}{\catname{LieGrpd}}
\newcommand{\sTop}{\catname{sTop}}
\newcommand{\Diff}{\catname{Diff}}
\newcommand{\sSet}{\catname{sSet}}
\newcommand{\Poset}{\catname{Poset}}

\renewcommand{\Pr}{\textup{Pr}}
\newcommand{\PrL}{\textup{Pr}^{\textup{L}}}
\newcommand{\PrR}{\textup{Pr}^{\textup{R}}}

\def\ulSp{\smash{\ul{{\setbox0=\hbox{\Sp}\dp0=1pt \ht0=0pt \box0\relax}}}}
\def\ulSpc{\smash{\ul{{\setbox0=\hbox{\Spc}\dp0=1pt \ht0=0pt \box0\relax}}}}
\def\ulOrbSp{\smash{\ul{{\setbox0=\hbox{\Orb\Sp}\dp0=1pt \ht0=0pt \box0\relax}}}}
\def\ulOrbSpc{\smash{\ul{{\setbox0=\hbox{\Orb\Spc}\dp0=1pt \ht0=0pt \box0\relax}}}}

%---------------------------------------
% Miscellaneous
%---------------------------------------

\newcommand{\h}{\textup{h}}
\newcommand{\CW}{\mathrm{CW}}
\DeclareMathOperator{\hocolim}{hocolim}
\DeclareMathOperator{\holim}{holim}
\DeclareMathOperator{\Exc}{Exc}
\newcommand{\s}{\textup{s}}
\DeclareMathOperator{\Sing}{Sing}
\DeclareMathOperator{\Kan}{Kan}
\newcommand{\bKan}{\mathbf{Kan}}
\newcommand{\bqCat}{\mathbf{qCat}}
\DeclareMathOperator{\Ch}{\catname{Ch}}
\DeclareMathOperator{\inj}{inj}
\DeclareMathOperator{\all}{all}
\DeclareMathOperator{\isotext}{iso}
\DeclareMathOperator{\Cut}{Cut}
\DeclareMathOperator{\LCut}{LCut}
\DeclareMathOperator{\BCut}{BCut}
\DeclareMathOperator{\Red}{Red}
\DeclareMathOperator{\Segal}{Segal}
\DeclareMathOperator{\Seg}{Seg}
\DeclareMathOperator{\CSeg}{CSeg}
\DeclareMathOperator{\add}{add}
\DeclareMathOperator{\sadd}{sadd}
\newcommand{\Piinfty}[1]{\Pi_{\infty}(#1)}
\DeclareMathOperator{\oplax}{oplax}

\newcommand{\qednow}{\pushQED{\qed}\qedhere\popQED} % Use this 

\newcommand{\quadtext}[1]{\quad\textrm{#1}\quad}
\newcommand{\qquadtext}[1]{\qquad\textrm{#1}\qquad}
\newcommand{\qin}{\quad\in\quad}
\newcommand{\ul}[1]{\underline{#1}}
\newcommand{\abs}[1]{\lvert #1 \rvert}
\newcommand{\geom}[1]{\| #1 \|}
\newcommand{\floor}[1]{\lfloor #1 \rfloor}

\newcommand{\catop}{^{\textup{op}}}
\let\op\relax
\DeclareMathOperator{\op}{op}

\DeclareMathOperator{\unit}{\mathds{1}}    % unit for \otimes
\newcommand{\blank}{{\textup{ --}}}
\newcommand{\ceil}[1]{\lceil #1 \rceil}
\newcommand{\lra}[1]{\langle #1 \rangle}

\newcommand{\iso}{\xrightarrow{\;\smash{\raisebox{-0.5ex}{\ensuremath{\scriptstyle\sim}}}\;}}
% \newcommand{\mapsfrom}{\mathrel{\reflectbox{$\mapsto$}}}
\newcommand{\xrightrightarrows}[2]{\begin{tikzcd}[ampersand replacement=\&, cramped] \vphantom{x}\rar[shift left=0.45ex, "#1"] \rar[shift right=0.45ex, swap, "#2"]\& \vphantom{x} \end{tikzcd}}
\newcommand{\xadjto}[2]{\begin{tikzcd}[ampersand replacement=\&, cramped] \vphantom{x}\rar[shift left=0.45ex, "#1"]\& \lar[shift left=0.45ex, "#2"] \vphantom{x} \end{tikzcd}}

%---------------------------------------
% Alphabets
%---------------------------------------
\let\S\relax
\let\P\relax
\DeclareMathOperator{\C}{\mathbb C}
\DeclareMathOperator{\F}{\mathbb F}
\DeclareMathOperator{\N}{\mathbb N}
\DeclareMathOperator{\P}{\mathbb P}
\DeclareMathOperator{\Q}{\mathbb Q}
\DeclareMathOperator{\R}{\mathbb R}
\DeclareMathOperator{\S}{\mathbb S}
\DeclareMathOperator{\Z}{\mathbb Z}
\DeclareMathOperator{\D}{\mathrm{D}}
\renewcommand{\phi}{\varphi}
\renewcommand{\epsilon}{\varepsilon}
\renewcommand{\aa}{\mathfrak{a}}
\newcommand{\pp}{\mathfrak{p}}
\newcommand{\fa}{\mathfrak{a}}
\newcommand{\fb}{\mathfrak{b}}
\newcommand{\fc}{\mathfrak{c}}
\newcommand{\fm}{\mathfrak{m}}
\newcommand{\mm}{\mathfrak{m}}
\newcommand{\fn}{\mathfrak{n}}
\newcommand{\fp}{\mathfrak{p}}

% Caligraphic font
\newcommand{\Aa}{\mathcal{A}}
\newcommand{\Bb}{\mathcal{B}}
\newcommand{\Cc}{\mathcal{C}}
\newcommand{\Dd}{\mathcal{D}}
\newcommand{\Ee}{\mathcal{E}}
\newcommand{\Ff}{\mathcal{F}}
\newcommand{\Gg}{\mathcal{G}}
\newcommand{\Hh}{\mathcal{H}}
\newcommand{\Ii}{\mathcal{I}}
\newcommand{\Jj}{\mathcal{J}}
\newcommand{\Kk}{\mathcal{K}}
\newcommand{\Ll}{\mathcal{L}}
\newcommand{\Mm}{\mathcal{M}}
\newcommand{\Nn}{\mathcal{N}}
\newcommand{\Oo}{\mathcal{O}}
\newcommand{\Pp}{\mathcal{P}}
\newcommand{\Qq}{\mathcal{Q}}
\newcommand{\Rr}{\mathcal{R}}
\newcommand{\Ss}{\mathcal{S}}
\newcommand{\Tt}{\mathcal{T}}
\newcommand{\Uu}{\mathcal{U}}
\newcommand{\Vv}{\mathcal{V}}
\newcommand{\Ww}{\mathcal{W}}
\newcommand{\Xx}{\mathcal{X}}
\newcommand{\Yy}{\mathcal{Y}}
\newcommand{\Zz}{\mathcal{Z}}

% Bold font
\newcommand{\bA}{\mathbf{A}}
\newcommand{\bB}{\mathbf{B}}
\newcommand{\bC}{\mathbf{C}}
\newcommand{\bD}{\mathbf{D}}
\newcommand{\bE}{\mathbf{E}}
\newcommand{\bF}{\mathbf{F}}
\newcommand{\bG}{\mathbf{G}}
\newcommand{\bH}{\mathbf{H}}
\newcommand{\bI}{\mathbf{I}}
\newcommand{\bJ}{\mathbf{J}}
\newcommand{\bK}{\mathbf{K}}
\newcommand{\bL}{\mathbf{L}}
\newcommand{\bM}{\mathbf{M}}
\newcommand{\bN}{\mathbf{N}}
\newcommand{\bO}{\mathbf{O}}
\newcommand{\bP}{\mathbf{P}}
\newcommand{\bQ}{\mathbf{Q}}
\newcommand{\bR}{\mathbf{R}}
\newcommand{\bS}{\mathbf{S}}
\newcommand{\bT}{\mathbf{T}}
\newcommand{\bU}{\mathbf{U}}
\newcommand{\bV}{\mathbf{V}}
\newcommand{\bW}{\mathbf{W}}
\newcommand{\bX}{\mathbf{X}}
\newcommand{\bY}{\mathbf{Y}}
\newcommand{\bZ}{\mathbf{Z}}

% Blackboard bold
\newcommand{\bbA}{\mathbb{A}}
\newcommand{\bbB}{\mathbb{B}}
\newcommand{\bbC}{\mathbb{C}}
\newcommand{\bbD}{\mathbb{D}}
\newcommand{\bbE}{\mathbb{E}}
\newcommand{\bbF}{\mathbb{F}}
\newcommand{\bbG}{\mathbb{G}}
\newcommand{\bbH}{\mathbb{H}}
\newcommand{\bbI}{\mathbb{I}}
\newcommand{\bbJ}{\mathbb{J}}
\newcommand{\bbK}{\mathbb{K}}
\newcommand{\bbL}{\mathbb{L}}
\newcommand{\bbM}{\mathbb{M}}
\newcommand{\bbN}{\mathbb{N}}
\newcommand{\bbO}{\mathbb{O}}
\newcommand{\bbP}{\mathbb{P}}
\newcommand{\bbQ}{\mathbb{Q}}
\newcommand{\bbR}{\mathbb{R}}
\newcommand{\bbS}{\mathbb{S}}
\newcommand{\bbT}{\mathbb{T}}
\newcommand{\bbU}{\mathbb{U}}
\newcommand{\bbV}{\mathbb{V}}
\newcommand{\bbW}{\mathbb{W}}
\newcommand{\bbX}{\mathbb{X}}
\newcommand{\bbY}{\mathbb{Y}}
\newcommand{\bbZ}{\mathbb{Z}}

%---------------------------------------
% Environments
%---------------------------------------

% Theorem-style environments
\theoremstyle{plain}
\newtheorem{theorem}{Theorem}[chapter]

\newtheorem{conjecture}[theorem]{Conjecture}
\newtheorem{corollary}[theorem]{Corollary}
\newtheorem{lemma}[theorem]{Lemma}
\newtheorem{proposition}[theorem]{Proposition}
\newtheorem{thmx}{Theorem}
\renewcommand{\thethmx}{\Alph{thmx}}

% Definition-style environments
\theoremstyle{definition}
\newtheorem{axiom}{Axiom}
\newtheorem{claim}[theorem]{Claim}
\newtheorem{construction}[theorem]{Construction}
\newtheorem{convention}[theorem]{Convention}
\newtheorem{definition}[theorem]{Definition}
\newtheorem{discussion}[theorem]{Discussion}
\newtheorem{example}[theorem]{Example}
\newtheorem{exercise}[theorem]{Exercise}
\newtheorem{fact}[theorem]{Fact}
\newtheorem{notation}[theorem]{Notation}
\newtheorem{observation}[theorem]{Observation}
\newtheorem{question}[theorem]{Question}
\newtheorem{remark}[theorem]{Remark}
\newtheorem{terminology}[theorem]{Terminology}
\newtheorem{warning}[theorem]{Warning}

% Unnumbered environments
\newtheorem*{remark*}{Remark}
\newtheorem*{example*}{Example}
\newtheorem*{theorem*}{Theorem}
\newtheorem*{recognition*}{Recognition Principle for Infinite Loop Spaces}
\newtheorem*{claim*}{Claim}
\newtheorem*{convention*}{Convention}
\newtheorem*{definition*}{Definition}
\newtheorem*{exercise*}{Exercise}
\newtheorem*{fact*}{Fact}
\newtheorem*{question*}{Question}
\newtheorem*{warning*}{Warning}
\renewcommand{\theaxiom}{\Alph{axiom}}
\newtheorem{subaxiom}{Axiom}
\renewcommand{\thesubaxiom}{\Alph{axiom}.\arabic{subaxiom}}
\renewcommand{\theHsubaxiom}{\theaxiom.\thesubaxiom}  % Makes sure the hyperlinks remember in which axiom the subaxiom appears
\newtheorem{subaxiomprime}{Axiom}
\renewcommand{\thesubaxiomprime}{\Alph{axiom}.\arabic{subaxiomprime}'}
\renewcommand{\theHsubaxiomprime}{\theaxiom.\thesubaxiomprime}  % Makes sure the hyperlinks remember in which axiom the subaxiom appears

%---------------------------------------
% Margin notes
%---------------------------------------
\newcommand{\red}[1]{{\color{red}{\textbf #1}}}
\setlength{\parindent}{0cm}
\setlength{\parskip}{0.8ex}
\setlength\marginparwidth{57pt}
\newcommand{\notehelper}[3]{\textcolor{#3}{$\blacksquare$}\marginpar{\ifodd\thepage\raggedright\else\raggedleft\fi\color{#3}\tiny \textbf{#2:} #1}}
\newcommand\bastiaan[1]{\notehelper{#1}{Bastiaan}{blue}}
\newcommand\bastiaanlong[1]{{\itshape \color{blue} #1}}

%---------------------------------------
% Tikzcd
%---------------------------------------

\tikzcdset{pullback/.style = {phantom, "\lrcorner"{very near start}}}
\tikzcdset{pullbackdl/.style = {phantom, "{\reflectbox{$\lrcorner$}}"{very near start}}}
\tikzcdset{pushout/.style = {phantom, "{\rotatebox{180}{$\lrcorner$}}"{very near end}}}

\tikzset{curve/.style={settings={#1},to path={(\tikztostart)
			.. controls ($(\tikztostart)!\pv{pos}!(\tikztotarget)!\pv{height}!270:(\tikztotarget)$)
			and ($(\tikztostart)!1-\pv{pos}!(\tikztotarget)!\pv{height}!270:(\tikztotarget)$)
			.. (\tikztotarget)\tikztonodes}}}
\tikzcdset{scale cd/.style={every label/.append style={scale=#1},cells={nodes={scale=#1}}}}
\newcommand{\pullback}{\mathrel{\text{\tikz{
				\draw[line width=0.6pt] (0,0) -- (0.5em,0);
				\draw[line width=0.6pt] (0.5em,0) -- (0.5em,0.6em);
		}}\,}}
\newcommand{\pullbackto}{\mathrel{\text{\tikz{
				\draw[line width=0.6pt] (0,0) -- (0.5em,0);
				\draw[line width=0.6pt] (0.5em,0) -- (0.5em,0.6em);
		}}\kern0.5em\to}}
\newcommand{\pushout}{\mathrel{\text{\tikz{
				\draw[line width=0.4pt] (0,0) -- (0,0.5em);
				\draw[line width=0.4pt] (0,0.5em) -- (0.5em,0.5em);
		}}\,}}
\newcommand{\pushoutto}{\mathrel{\text{\tikz{
				\draw[line width=0.4pt] (0,0) -- (0,0.5em);
				\draw[line width=0.4pt] (0,0.5em) -- (0.5em,0.5em);
		}}\kern0.5em\to}}

%---------------------------------------
% Other
%---------------------------------------

\newcommand{\quot}{/\!\!/}

% Use this when you want to write a \colon but the arrow goes from right to left, as in $Y \leftarrow X\noloc f$.
\newcommand\noloc{%
	\nobreak
	\mspace{6mu plus 1mu}
	{:}
	\nonscript\mkern-\thinmuskip
	\mathpunct{}
	\mspace{2mu}
}

\newcommand{\pvto}[2]{{%
		\kern1pt
		\big(
		\begin{matrix}
			\text{\tiny $#1$} \\[-7pt]
			\scriptstyle\downarrow \\[-7pt]
			\text{\tiny $#2$}
		\end{matrix}
		\big)
}}
\newcommand{\puto}[3]{{%
		\kern1pt
		\big(
		\begin{matrix}
			\text{\tiny $#1$ \hspace{-6pt}\phantom{$#3$}} \\[-9pt]
			\scriptstyle\downarrow \raisebox{2pt}{\text{\tiny $#3$}} \\[-9pt]
			\text{\tiny $#2$ \hspace{-6pt}\phantom{$#3$}}
		\end{matrix}
		\big)
}}
\newcommand{\apvto}[2]{\raisebox{0pt}[1.3\ht\strutbox][1.3\dp\strutbox]{$\smash{\pvto{#1}{#2}}$}}
\newcommand{\pwto}[2]{{%
		\kern1pt
		\big(
		\begin{matrix}
			\text{\tiny $#1$} \\[-5pt]
			\text{\tiny $#2$}
		\end{matrix}
		\big)
}}
\newcommand{\rbox}[3]{\raisebox{0pt}[#1\ht\strutbox][#2\dp\strutbox]{$#3$}}

%---------------------------------------
% Cleverref
%---------------------------------------
\usepackage[capitalise]{cleveref}  % Smart cross-referencing

\crefname{axiom}{axiom}{axioms}
\crefname{claim}{claim}{claims}
\crefname{conjecture}{conjecture}{conjectures}
\crefname{construction}{construction}{constructions}
\crefname{convention}{convention}{conventions}
\crefname{corollary}{corollary}{corollaries}
\crefname{definition}{definition}{definitions}
\crefname{discussion}{discussion}{discussions}
\crefname{example}{example}{examples}
\crefname{exercise}{exercise}{exercises}
\crefname{fact}{fact}{facts}
\crefname{lemma}{lemma}{lemmas}
\crefname{notation}{notation}{notations}
\crefname{observation}{observation}{observations}
\crefname{proposition}{proposition}{propositions}
\crefname{question}{question}{questions}
\crefname{remark}{remark}{remarks}
\crefname{subaxiom}{axiom}{axioms}
\crefname{terminology}{terminology}{terminologies}
\crefname{theorem}{theorem}{theorems}
\crefname{warning}{warning}{warnings}

\Crefname{axiom}{Axiom}{Axioms}
\Crefname{claim}{Claim}{Claims}
\Crefname{conjecture}{Conjecture}{Conjectures}
\Crefname{construction}{Construction}{Constructions}
\Crefname{convention}{Convention}{Conventions}
\Crefname{corollary}{Corollary}{Corollaries}
\Crefname{definition}{Definition}{Definitions}
\Crefname{discussion}{Discussion}{Discussions}
\Crefname{example}{Example}{Examples}
\Crefname{exercise}{Exercise}{Exercises}
\Crefname{fact}{Fact}{Facts}
\Crefname{lemma}{Lemma}{Lemmas}
\Crefname{notation}{Notation}{Notations}
\Crefname{observation}{Observation}{Observations}
\Crefname{proposition}{Proposition}{Propositions}
\Crefname{question}{Question}{Questions}
\Crefname{remark}{Remark}{Remarks}
\Crefname{subaxiom}{Axiom}{Axioms}
\Crefname{terminology}{Terminology}{Terminologies}
\Crefname{theorem}{Theorem}{Theorems}
\Crefname{warning}{Warning}{Warnings}

% Appendix chapters are labelled \label[appendix]{...} so that \Cref calls them
% ``Appendix'' rather than ``Chapter'' (main-matter chapters keep ``Chapter'').
\crefname{appendix}{appendix}{appendices}
\Crefname{appendix}{Appendix}{Appendices}
